% ============================================================
% M3 S2 导学件·波1 臂6 片件 —— 课时11 椭圆的标准方程（2.5.1）＝批C 四片之二（母版照抄制）
% 生成：M3 成卷轮 S2 波1 臂6｜2026-09-14｜编译 cwd＝本目录（中文路径禁 \input 绝对路径）
%   xelatex main-true.tex ×2 → 含详解印本；xelatex main-false.tex ×2 → 纯题版（单源双本）
% 输入链（全只读）：题面库 课时11-2.5.1椭圆的标准方程.md（题面逐字装配）＋ -答案侧.md（值逐字回嵌）
%   ＋ 定稿/批C-课时11-2.5.1椭圆的标准方程.md（详解回嵌源，§9 全文区）；toolchain＝qp-m3.sty
%   （钉后版 md5 7c3930362be8a0a2bdf21bbf8ac16573，本地挂载副本，破损即停）。
% 母版体例：详见 成卷/导学件/课时01-坐标法/母版体例说明.md（结构骨架/宏用法/门谱跑法/常见坑）。
% 键账：件内 ansblock 键集 ≡ 题面库片键集 ≡ manifest 键序（21 键，零缺漏零多余）；
%   预习位零键零答案（口令④裁定前口径）；印面号 1—21＝探究 1—16＋课堂评价 17—21，
%   键↔号映射见 值台账-课时11.json。
% 括线判据（承重墙禁放宽）：估高＞8 行（chars/23 栏宽口径）→ 括线模，逐键读数入值台账；
%   其余灰底。本片括线键读数见值台账（门谱实测）。
% 观察注落实（军师令·补产批CA 核账 §C.5 随片下发）：难2＝题面照录在印＋「装配替换建议
%   （S4可换问消前引）」入值台账；简10＝「离心率」前引同难2，替换建议入值台账；
%   简3＝详解承定稿直验口径（两位置简验）照贴；中3＝白名单题面照录在印；
%   中6（拓16席）/拓2＝非本件键域，本件落实数 0；逐条读数见 值台账-课时11.json「观察注落实」。
% 门谱读数：见过程件 工作区/_tmpM3S2波1臂6_0914/（编译 log、门谱输出、键表）。
% ============================================================
\documentclass[fontset=none]{ctexart}
\usepackage{qp-m3}
% —— 印面逐字符号路由补钉（探针实证 0914：书宋缺 U+2212，▱ 诸字库皆缺→NSC 子块；
%     禁改 sty 本体保 md5；无 ▱/− 用件的片可整块照抄，零副作用） ——
\xeCJKDeclareCharClass{Default}{"2212}%
\setCJKfamilyfont{nscblk}[Path=C:/提示词/工作区/字替对照-0909/variantF/fonts/,BoldFont=NSC-w600.ttf]{NSC-w500.ttf}%
\xeCJKDeclareSubCJKBlock{nscblk}{"25B1}%
\newcommand{\pxparallelogram}{{\CJKfamily{nscblk}▱}}%
% —— 断行微调（纯题档/详解档三零门：CJK 胶收缩 0.01→0.05em；只动伸缩分量不动自然宽度） ——
\xeCJKsetup{CJKglue={\hskip 0.10em plus 0.08em minus 0.05em}}%
% —— 页脚件名（qp-headfoot 复用入口） ——
\renewcommand{\qpzhangming}{第二章\quad 平面解析几何}%
\renewcommand{\qpjianming}{导学件}%

\begin{document}
%装配轮S6三本跨本连续页码（G7方案B）setcounter 注入位
\setcounter{page}{67}

\zhangtitle{第二章\quad 平面解析几何}
\jietitle{2.5.1 椭圆的标准方程}
\mubiaoline
\mubiaomu{1}{理解椭圆的定义，能用定义辨析动点轨迹（椭圆、线段、无轨迹三段），掌握定值与焦距的大小关系；}
\mubiaomu{2}{掌握椭圆的标准方程（焦点在\(x\)轴、\(y\)轴两形），会用定义法、待定系数法依条件求椭圆的标准方程，能由方程辨识焦点位置与基本量\(a\)、\(b\)、\(c\)；}
\mubiaomu{3}{会用椭圆定义处理焦点三角形（周长、面积、焦半径和与积的最值）及向量构形问题，体会“回到定义去”的解析几何思想.}

\vspace{1.7mm}
\begin{multicols}{2}
\emergencystretch=1em

% ============================================================
% 课前预习（知识点导学填空；零键零答案——预习键位按检查口令④裁定前口径，
% \kongbai 空线＋\zhentib 判断题面形，两档等形不泄答）
% ============================================================
\huaxing{课}{前}{预}{习}{知识导学\quad 素养初识}

\zsd{一}{椭圆的定义}

\tiaomu{1}{平面内与两个定点\(F_1\)、\(F_2\)的距离的\kongbai{}等于常数（大于\(|F_1F_2|\)）的点的轨迹叫作\kongbai{}；两定点叫作椭圆的\kongbai{}，两焦点间的距离叫作椭圆的\kongbai{}．\zhuzhu{定义两要素：定值\(2a\)、定值大于焦距\(2c\)；\(2a=|F_1F_2|\)时轨迹为线段\(F_1F_2\)，\(2a<|F_1F_2|\)时无轨迹，辨析题就查这两处.}}

\zsd{二}{椭圆的标准方程}

\tiaomu{1}{焦点在\(x\)轴上：\(x^{2}/a^{2}+y^{2}/b^{2}=1\)（\(a>b>0\)），焦点\((\kongbai{},0)\)；焦点在\(y\)轴上：\(y^{2}/a^{2}+x^{2}/b^{2}=1\)，焦点\((0,\kongbai{})\)；三个基本量满足\kongbai{}．\zhuzhu{焦点位置看大分母：分母大者在哪条轴，焦点就在哪条轴；\(a\)恒为最大半轴，\(a^{2}=b^{2}+c^{2}\)（与双曲线差号，勿混）.}}

\zsd{三}{依条件求标准方程·定义的应用}

\tiaomu{1}{由定义直读：焦点已知则\(c\)可读，距离和给出\(2a\)后\(b^{2}=a^{2}-c^{2}\)；双条件（距离和＋通径长/过点等）用\kongbai{}法解出\(a^{2}\)、\(b^{2}\)．\zhuzhu{焦点三角形常用拼合：过焦点的弦与两焦半径凑成周长\(4a\)；焦半径和恒为\(2a\)；焦半径积的最值用均值不等式，等号在短轴端点.}}

\tiaomu{2}{含参方程\(x^{2}/m+y^{2}/n=1\)表示椭圆要求\(m\)、\(n\)同正；焦点位置比较\(m\)、\(n\)的\kongbai{}．\zhuzhu{分母含参先保证均为正才谈位置；“焦点在\(y\)轴”即\(y^{2}\)分母较大.}}

\zhenhead{判断正误(正确的打\gou{},错误的打\cha{})}

\zhentib{(1)}{平面内到两个定点距离之和为常数（大于零）的点的轨迹一定是椭圆．}

\zhentib{(2)}{椭圆\(x^{2}/9+y^{2}/5=1\)的焦点在\(y\)轴上．}

\zhentib{(3)}{椭圆上任意一点到其两焦点的距离之和都等于\(2a\)．}

\liubai[9mm]{此处书写}

% ============================================================
% 课中探究（考点探究〔例+变式〕＝练/拓槽 16 键；题后紧跟答案＝ansblock 内嵌，
% 值照答案侧逐字，详解承定稿 §9 回嵌）
% ============================================================
\huaxing{课}{中}{探}{究}{考点探究\quad 素养小结}

% ---------- 探究点一 椭圆的定义与轨迹辨析 ----------
\tjdnr{一}{椭圆的定义与轨迹辨析}{例\textbf{1}}{简单(知识点一)}{已知定点\(F_1(-4,0)\)、\(F_2(4,0)\)和动点\(M(x,y)\)．(1)再从条件①、条件②这两个条件中选择一个作为已知，求：动点\(M\)的轨迹及其方程．条件①：\(|MF_1|+|MF_2|=12\)；条件②：\(|MF_1|+|MF_2|=8\)．(2)\(|MF_1|+|MF_2|=2a\)（\(a>0\)），求：动点\(M\)的轨迹及其方程．}
\xiexwei{12mm}

{\ansblockgrayfalse
\begin{ansblock}[2章-练-课时11-简7]
% ans:2章-练-课时11-简7
\ansitem{1}{(1)选①：椭圆x²/36+y²/20=1；选②：线段F₁F₂，方程y=0（−4≤x≤4）；(2)0＜a＜4无轨迹；a=4线段y=0（−4≤x≤4）；a＞4椭圆x²/a²+y²/(a²−16)=1}
\ansnote{详解}{(1)选①：\(12>|F_1F_2|=8\)，轨迹是以\(F_1\)、\(F_2\)为焦点的椭圆，\(c=4\)、\(a=6\)、\(b^{2}=36-16=20\)：\(x^{2}/36+y^{2}/20=1\)．选②：\(8=|F_1F_2|\)，距离和等于焦距，\(M\)只能在线段\(F_1F_2\)上（三角形不等式取等），轨迹为线段，方程\(y=0(-4\leqslant x\leqslant 4)\)．(2)按\(2a\)与\(|F_1F_2|=8\)比较三段：\(0<a<4\)时距离和小于焦距，无点满足，无轨迹；\(a=4\)同上为线段；\(a>4\)为椭圆，\(c=4\)、\(b^{2}=a^{2}-16\)：\(x^{2}/a^{2}+y^{2}/(a^{2}-16)=1\)．综上如答案．}
\end{ansblock}}

\liB{变式\textbf{2}}{简单(知识点一)}{}{已知\(F_1(-2,0)\)、\(F_2(2,0)\)，动点\(P\)满足\(|PF_1|+|PF_2|=6\)，则\(P\)的轨迹方程为\kongbai{}．}
\xiexwei{10mm}

\begin{ansblock}[2章-练-课时11-简8]
% ans:2章-练-课时11-简8
\ansitem{2}{x²/9＋y²/5＝1}
\ansnote{详解}{\(|F_1F_2|=4<6\)，由椭圆定义：\(2a=6\)，\(a=3\)；\(c=2\)；\(b^{2}=a^{2}-c^{2}=9-4=5\)．焦点在\(x\)轴，方程\(x^{2}/9+y^{2}/5=1\)．验算：取顶点\((3,0)\)：\(|PF_1|+|PF_2|=5+1=6\)成立；取\((0,\sqrt{5})\)：\(2\times\sqrt{4+5}=6\)成立．辨析：若和\(=4\)为线段\(F_1F_2\)，若和\(<4\)无轨迹（定义辨析位已讲，本题为\(>4\)基础位）．}
\end{ansblock}

\liB{变式\textbf{3}}{简单(知识点一)}{}{已知点\(A(4,0)\)，\(B(2,2)\)是椭圆\(x^{2}/25+y^{2}/9=1\)内的两个点，\(M\)是椭圆上的动点，则\(|MA|+|MB|\)的最大值为\kongbai{}．}
\xiexwei{10mm}

\begin{ansblock}[2章-练-课时11-简2]
% ans:2章-练-课时11-简2
\ansitem{3}{10+2√10}
\ansnote{详解}{\(a=5\)、\(b=3\)、\(c=4\)，\(A(4,0)\)恰为右焦点，记左焦点\(C(-4,0)\)．由定义\(|MA|=2a-|MC|=10-|MC|\)，故\(|MA|+|MB|=10+(|MB|-|MC|)\leqslant 10+|BC|=10+\sqrt{(2+4)^{2}+(2-0)^{2}}=10+\sqrt{40}=10+2\sqrt{10}\)．等号当\(M\)在射线\(BC\)（自\(B\)过\(C\)延长方向）与椭圆的交点处取得（\(C\)在椭圆内，射线必与椭圆交于一点），可取．最大值\(10+2\sqrt{10}\)．}
\end{ansblock}

\liB{变式\textbf{4}}{简单(知识点一)}{}{设\(F_1\)、\(F_2\)是椭圆\(x^{2}/16+y^{2}/4=1\)的左、右焦点，过\(F_1\)的直线\(l\)交椭圆于\(A\)、\(B\)两点，则\(|AF_2|+|BF_2|\)的最大值为\nobreak(\kongwei)}
\bindopt {\fontsize{10.5pt}{\qpTJgLeadLiti}\selectfont \makebox[\dimexpr0.5\linewidth-1em\relax][l]{A．\(14\)}\makebox[\dimexpr0.5\linewidth-1em\relax][l]{B．\(13\)}\par}

\bindopt {\fontsize{10.5pt}{\qpTJgLeadLiti}\selectfont \makebox[\dimexpr0.5\linewidth-1em\relax][l]{C．\(12\)}\makebox[\dimexpr0.5\linewidth-1em\relax][l]{D．\(10\)}\par}

{\ansblockgrayfalse
\begin{ansblock}[2章-练-课时11-简3]
% ans:2章-练-课时11-简3
\ansitem{4}{A}
\ansnote{详解}{\(a=4\)、\(b=2\)．定义：\(|AF_1|+|AF_2|=|BF_1|+|BF_2|=8\)，两式相加并注意\(|AB|=|AF_1|+|BF_1|\)（\(A\)、\(F_1\)、\(B\)共线）：\(\triangle ABF_2\)周长\(|AF_2|+|BF_2|+|AB|=16=4a\)．故\(|AF_2|+|BF_2|=16-|AB|\)，要最大须过焦点的弦\(|AB|\)最小．竖直位置：\(l\)：\(x=-c=-2\sqrt{3}\)，\(y^{2}=4(1-12/16)=1\)，\(|AB|=2\)（通径\(2b^{2}/a=2\)）；水平位置：长轴弦\(|AB|=2a=8\)．通径为最短焦点弦（一般证明需弦长联立，2.5.3后；此处按两位置比较简验）．最大值\(=16-2=14\)，故选：A．}
\end{ansblock}}

\liB{变式\textbf{5}}{简单(知识点一)}{}{若\(\triangle ABC\)中，\(|BC|=m\)，\(|AB|+|AC|=n\)（\(0<m<n\)，且\(m\)、\(n\)为定值），则\(\triangle ABC\)面积的最大值为\kongbai{}．}
\xiexwei{12mm}

{\ansblockgrayfalse
\begin{ansblock}[2章-练-课时11-简4]
% ans:2章-练-课时11-简4
\ansitem{5}{(m/4)√(n²−m²)}
\ansnote{详解}{\(n>m\)：\(A\)到两定点\(B\)、\(C\)距离和为常数\(n\)且大于\(|BC|\)，由定义\(A\)在以\(B\)、\(C\)为焦点的椭圆上（\(2a=n\)、\(2c=m\)，\(A\)不在直线\(BC\)上故去长轴端点）．以\(BC\)为底，底长\(m\)固定，面积最大即高最大；椭圆上点到焦点连线所在直线的最大距离为短半轴长\(b=\sqrt{a^{2}-c^{2}}=\sqrt{(n/2)^{2}-(m/2)^{2}}=\frac{1}{2}\sqrt{n^{2}-m^{2}}\)（短轴端点取得，该点不与\(B\)、\(C\)共线，合法）．\(S_{\max}=\frac{1}{2}\cdot m\cdot\frac{1}{2}\sqrt{n^{2}-m^{2}}=(m/4)\sqrt{n^{2}-m^{2}}\)．}
\end{ansblock}}

\xiaojie{①识别：给定距离和与定点的轨迹辨析、距离和的最值（接焦点转化）、定义存在性判断，判归本题型.}
\bindp {\kaishu ②操作：先比定值\(2a\)与焦距\(|F_1F_2|\)定三段（大——椭圆；等——线段；小——无轨迹）；椭圆内定点接焦点，“到焦点距离”与“\(2a-\)到另一焦点距离”互换后再用三角形不等式.}
\bindp {\kaishu ③收束：等号可取性要交代（交点/端点在不在轨迹上）；线段轨迹的方程要带范围限定（如\(y=0(-4\leqslant x\leqslant 4)\)）.}

% ---------- 探究点二 依条件求标准方程 ----------
\tjdnr{二}{依条件求标准方程}{例\textbf{6}}{简单(知识点二)}{已知椭圆\(C\)：\(x^{2}/a^{2}+y^{2}/b^{2}=1\)（\(a>b>0\)）的左、右焦点分别是\(F_1\)、\(F_2\)，椭圆\(C\)上任意一点到\(F_1\)、\(F_2\)的距离之和为\(4\sqrt{3}\)，过焦点\(F_2\)且垂直于\(x\)轴的直线交椭圆于\(A\)、\(B\)两点，若线段\(AB\)的长度为\(4\sqrt{3}/3\)，则椭圆\(C\)的方程为\kongbai{}．}
\xiexwei{12mm}

{\ansblockgrayfalse
\begin{ansblock}[2章-练-课时11-简1]
% ans:2章-练-课时11-简1
\ansitem{6}{x²/12+y²/4=1}
\ansnote{详解}{定义：\(2a=4\sqrt{3}\)，即\(a=2\sqrt{3}\)、\(a^{2}=12\)．设\(A(c,y_0)\)：\(c^{2}/a^{2}+y_0^{2}/b^{2}=1\)，即\(y_0^{2}=b^{2}(1-c^{2}/a^{2})=b^{4}/a^{2}\)，\(|y_0|=b^{2}/a\)；通径\(|AB|=2b^{2}/a=4\sqrt{3}/3\)，即\(b^{2}=a\cdot(2\sqrt{3})/3=2\sqrt{3}\cdot(2\sqrt{3})/3=4\)．椭圆方程\(x^{2}/12+y^{2}/4=1\)．验算：\(c^{2}=12-4=8\)，\(x=c=2\sqrt{2}\)代入得\(y^{2}=4(1-8/12)=4/3\)，\(|AB|=2\times(2/\sqrt{3})=4\sqrt{3}/3\)成立；顶点\((\sqrt{12},0)\)到两焦点\((\pm 2\sqrt{2},0)\)距离和\(=(2\sqrt{2}+\sqrt{12})+(\sqrt{12}-2\sqrt{2})=2\sqrt{12}=4\sqrt{3}\)成立．方程正确．}
\end{ansblock}}

\liB{变式\textbf{7}}{简单(知识点二)}{}{方程\(x^{2}/(m-2)+y^{2}/(6-m)=1\)表示焦点在\(y\)轴上的椭圆，求实数\(m\)的取值范围．}
\xiexwei{10mm}

\begin{ansblock}[2章-练-课时11-简9]
% ans:2章-练-课时11-简9
\ansitem{7}{2＜m＜4}
\ansnote{详解}{椭圆成立需\(m-2>0\)且\(6-m>0\)，即\(2<m<6\)；焦点在\(y\)轴需\(y^{2}\)分母较大：\(6-m>m-2\)，即\(m<4\)．综上\(2<m<4\)．验算：\(m=3\)：\(x^{2}/1+y^{2}/3=1\)，\(a^{2}=3>b^{2}=1\)，焦点\((0,\pm\sqrt{2})\)在\(y\)轴成立；\(m=5\)：\(x^{2}/3+y^{2}/1=1\)焦点在\(x\)轴（排除端外值实证）成立．}
\end{ansblock}

\liB{变式\textbf{8}}{简单(知识点二)}{}{求中心在原点、对称轴为坐标轴、焦点在\(x\)轴上，且过点\(P(2,1)\)、离心率\(e=1/2\)的椭圆方程．}
\xiexwei{12mm}

{\ansblockgrayfalse
\begin{ansblock}[2章-练-课时11-简10]
% ans:2章-练-课时11-简10
\ansitem{8}{x²/(16/3)＋y²/4＝1}
\ansnote{详解}{\(e=c/a=1/2\)，即\(a=2c\)，\(b^{2}=a^{2}-c^{2}=3c^{2}\)．设\(x^{2}/a^{2}+y^{2}/b^{2}=1\)，过\(P\)：\(4/a^{2}+1/b^{2}=1\)．代\(a^{2}=4c^{2}\)、\(b^{2}=3c^{2}\)：\(4/(4c^{2})+1/(3c^{2})=1\)，即\(\frac{1}{c^{2}}+\frac{1}{3c^{2}}=1\)，得\(\frac{4}{3c^{2}}=1\)，\(c^{2}=4/3\)．故\(a^{2}=16/3\)，\(b^{2}=4\)．方程\(x^{2}/(16/3)+y^{2}/4=1\)．验算：\(P(2,1)\)：\(4/(16/3)+1/4=12/16+4/16=1\)成立；\(e\)：\(c=2/\sqrt{3}\)，\(a=4/\sqrt{3}\)，\(e=1/2\)成立．}
\end{ansblock}}

\liB{变式\textbf{9}}{简单(知识点二)}{}{已知点\(M\)是椭圆\(x^{2}/a^{2}+y^{2}/b^{2}=1\)（\(a>b>0\)）上任意一点，两个焦点分别为\(F_1\)、\(F_2\)，若\(|MF_1|\cdot|MF_2|\)的最大值为8，则\(a\)的值为\nobreak(\kongwei)}
\bindopt {\fontsize{10.5pt}{\qpTJgLeadLiti}\selectfont \makebox[\dimexpr0.5\linewidth-1em\relax][l]{A．\(8\)}\makebox[\dimexpr0.5\linewidth-1em\relax][l]{B．\(4\)}\par}

\bindopt {\fontsize{10.5pt}{\qpTJgLeadLiti}\selectfont \makebox[\dimexpr0.5\linewidth-1em\relax][l]{C．\(2\sqrt{2}\)}\makebox[\dimexpr0.5\linewidth-1em\relax][l]{D．\(2\)}\par}

\begin{ansblock}[2章-练-课时11-简5]
% ans:2章-练-课时11-简5
\ansitem{9}{C}
\ansnote{详解}{定义\(|MF_1|+|MF_2|=2a\)，均值不等式\(|MF_1|\cdot|MF_2|\leqslant((|MF_1|+|MF_2|)/2)^{2}=a^{2}\)，等号当\(|MF_1|=|MF_2|\)（\(M\)为短轴端点，可取）．最大值\(a^{2}=8\)，即\(a=2\sqrt{2}\)（\(a>0\)）．故选：C．}
\end{ansblock}

\xiaojie{①识别：由距离和/通径/过点/焦点位置等条件反求椭圆方程，判归本题型.}
\bindp {\kaishu ②操作：定义直读\(a\)、\(c\)则\(b^{2}=a^{2}-c^{2}\)；否则待定系数设标准形列方程组；含参方程先保分母同正，再看大分母定焦点位置.}
\bindp {\kaishu ③收束：焦点位置认准大分母一侧，\(a\)恒最大；求出方程回代条件验算一遍防设形错位.}

% ---------- 探究点三 焦点三角形与向量构形 ----------
\tjdnr{三}{焦点三角形与向量构形}{例\textbf{10}}{简单(知识点三)}{\duoxuan\hspace{0.6em}点\(F_1\)、\(F_2\)为椭圆\(C\)的两个焦点，椭圆\(C\)上存在点\(P\)，使得\(\angle F_1PF_2=90^{\circ}\)，则椭圆\(C\)的方程可以是\nobreak(\kongwei)}
\bindopt {\fontsize{10.5pt}{\qpTJgLeadLiti}\selectfont \makebox[\dimexpr0.5\linewidth-1em\relax][l]{A．\(x^{2}/25+y^{2}/9=1\)}\makebox[\dimexpr0.5\linewidth-1em\relax][l]{B．\(x^{2}/25+y^{2}/16=1\)}\par}

\bindopt {\fontsize{10.5pt}{\qpTJgLeadLiti}\selectfont \makebox[\dimexpr0.5\linewidth-1em\relax][l]{C．\(x^{2}/18+y^{2}/9=1\)}\makebox[\dimexpr0.5\linewidth-1em\relax][l]{D．\(x^{2}/16+y^{2}/8=1\)}\par}

{\ansblockgrayfalse
\begin{ansblock}[2章-练-课时11-简6]
% ans:2章-练-课时11-简6
\ansitem{10}{ACD}
\ansnote{详解}{\(\angle F_1PF_2=90^{\circ}\)当且仅当\(P\)在以\(F_1F_2\)为直径的圆\(x^{2}+y^{2}=c^{2}\)上（圆周角定理，\(P\neq F_1\)、\(F_2\)自动满足）．设\(P(x,y)\)在椭圆上，则\(|OP|^{2}=x^{2}+y^{2}=x^{2}+b^{2}(1-x^{2}/a^{2})=b^{2}+x^{2}(1-b^{2}/a^{2})\)，\(x^{2}\in[0,a^{2}]\)且\(1-b^{2}/a^{2}>0\)，故\(|OP|^{2}\in[b^{2},a^{2}]\)，即椭圆上点到原点距离介于\(b\)与\(a\)之间．圆与椭圆有公共点当且仅当\(b^{2}\leqslant c^{2}\leqslant a^{2}\)，右半恒真，条件即\(b\leqslant c\)，即\(b^{2}\leqslant a^{2}-b^{2}\)，即\(a^{2}\geqslant 2b^{2}\)（存在性与上顶点处张角\(\angle F_1BF_2\geqslant 90^{\circ}\)同述）．逐项：A \(25\geqslant 18\)成立；B \(25\geqslant 32\)不成立；C \(18\geqslant 18\)成立；D \(16\geqslant 16\)成立．故选：ACD．}
\end{ansblock}}

\liB{变式\textbf{11}}{中档(知识点三)}{}{已知椭圆\(C\)：\(x^{2}/a^{2}+y^{2}/b^{2}=1\)（\(a>b>0\)）的焦点为\(F_1\)、\(F_2\)，若椭圆\(C\)上存在一点\(P\)，使得向量\(PF_1\cdot\)向量\(PF_2=0\)，且\(\triangle F_1PF_2\)的面积等于4，则实数\(b\)的值为\kongbai{}．}
\xiexwei{12mm}

{\ansblockgrayfalse
\begin{ansblock}[2章-练-课时11-中1]
% ans:2章-练-课时11-中1
\ansitem{11}{2}
\ansnote{详解}{向量\(PF_1\cdot\)向量\(PF_2=0\)即\(\angle F_1PF_2=90^{\circ}\)，则\(|PF_1|^{2}+|PF_2|^{2}=|F_1F_2|^{2}=4c^{2}\)．记\(m=|PF_1|\cdot|PF_2|\)：面积\(\frac{1}{2}m=4\)，即\(m=8\)．定义：\((|PF_1|+|PF_2|)^{2}=4a^{2}=|PF_1|^{2}+|PF_2|^{2}+2|PF_1||PF_2|=4c^{2}+16\)，即\(a^{2}-c^{2}=4\)，\(b^{2}=4\)，\(b>0\)得\(b=2\)．验算相容性：\(|PF_1|\)、\(|PF_2|\)是\(t^{2}-2at+8=0\)的正根，判别式\(4a^{2}-32\geqslant 0\)即\(a^{2}\geqslant 8\)（题设“存在\(P\)”即椭圆满足此，必要条件下\(b\)恒为2）．}
\end{ansblock}}

\liB{变式\textbf{12}}{中档(知识点三)}{}{经过椭圆\(C\)：\(x^{2}/2+y^{2}=1\)的中心作直线\(l\)，与椭圆\(C\)交于\(P\)、\(Q\)两点，设椭圆\(C\)的右焦点为\(F\)，已知\(\angle PFQ=2\pi/3\)，求\(\triangle PFQ\)的面积．}
\xiexwei{12mm}

{\ansblockgrayfalse
\begin{ansblock}[2章-练-课时11-中2]
% ans:2章-练-课时11-中2
\ansitem{12}{√3/3}
\ansnote{详解}{记左焦点\(F_1\)．\(O\)既是\(PQ\)中点（\(l\)过中心、椭圆关于\(O\)对称）又是\(FF_1\)中点，故四边形\(PF_1QF\)对角线互相平分，为平行四边形．于是\(S_{\triangle PFQ}=S_{\triangle PF_1F}\)（对角线\(F_1F\)所分另一半等积）．平行四边形邻角互补：顶点\(P\)处内角\(\angle F_1PF=\pi-\angle PFQ=\pi/3\)．在\(\triangle PF_1F\)中：\(a=\sqrt{2}\)、\(b=1\)、\(c=1\)，\(|PF_1|+|PF|=2\sqrt{2}\)，\(|F_1F|=2\)，\(\angle F_1PF=\pi/3\)．余弦定理：\(4=|PF_1|^{2}+|PF|^{2}-2|PF_1||PF|\cos(\pi/3)=(2\sqrt{2})^{2}-3|PF_1||PF|\)，即\(|PF_1||PF|=4/3\)．\(S_{\triangle PF_1F}=\frac{1}{2}\cdot\frac{4}{3}\cdot\sin(\pi/3)=\sqrt{3}/3\)，即\(\triangle PFQ\)的面积．验算：积\(4/3\leqslant((2\sqrt{2})/2)^{2}=4\)成立，两半径可正实存在．}
\end{ansblock}}

\liB{变式\textbf{13}}{中档(知识点三)}{}{在平面直角坐标系\(xOy\)中，已知\(\triangle ABC\)顶点\(A(-1,0)\)和\(C(1,0)\)，顶点\(B\)在椭圆\(x^{2}/4+y^{2}/3=1\)上，则\((\sin A+\sin C)/\sin B\)的值是\nobreak(\kongwei)}
\bindopt {\fontsize{10.5pt}{\qpTJgLeadLiti}\selectfont \makebox[\dimexpr0.5\linewidth-1em\relax][l]{A．\(0\)}\makebox[\dimexpr0.5\linewidth-1em\relax][l]{B．\(1\)}\par}

\bindopt {\fontsize{10.5pt}{\qpTJgLeadLiti}\selectfont \makebox[\dimexpr0.5\linewidth-1em\relax][l]{C．\(2\)}\makebox[\dimexpr0.5\linewidth-1em\relax][l]{D．不确定}\par}

\begin{ansblock}[2章-练-课时11-中3]
% ans:2章-练-课时11-中3
\ansitem{13}{C}
\ansnote{详解}{\(A\)、\(C\)恰为椭圆两焦点（\(c=1\)，\(a=2\)）．\(B\)在椭圆上且\(\triangle ABC\)非退化（\(B\)不为长轴端点，\(\sin B\neq 0\)）．正弦定理：\(|BC|=2R'\sin A\)、\(|AB|=2R'\sin C\)、\(|AC|=2R'\sin B\)（\(R'\)为\(\triangle ABC\)外接圆半径），故\((\sin A+\sin C)/\sin B=(|BC|+|AB|)/|AC|\)．由定义\(|AB|+|BC|=2a=4\)，\(|AC|=2c=2\)，比值\(=4/2=2\)，与\(B\)位置无关．故选：C．}
\end{ansblock}

\liB{变式\textbf{14}}{中档(知识点三)}{}{设\(F_1\)、\(F_2\)分别为椭圆\(x^{2}/3+y^{2}=1\)的左、右焦点，点\(A\)、\(B\)在椭圆上．若向量\(F_1A=5\cdot\)向量\(F_2B\)，则点\(A\)的坐标是\kongbai{}．}
\xiexwei{12mm}

{\ansblockgrayfalse
\begin{ansblock}[2章-练-课时11-中4]
% ans:2章-练-课时11-中4
\ansitem{14}{(0,±1)}
\ansnote{详解}{\(a^{2}=3\)、\(b^{2}=1\)、\(c^{2}=2\)：\(F_1(-\sqrt{2},0)\)、\(F_2(\sqrt{2},0)\)．设\(A(x_1,y_1)\)、\(B(x_2,y_2)\)．由向量等式：\((x_1+\sqrt{2},y_1)=5(x_2-\sqrt{2},y_2)\)，即\(x_2=(x_1+6\sqrt{2})/5\)、\(y_2=y_1/5\)．两点均在椭圆上：\(x_1^{2}/3+y_1^{2}=1\)且\((x_1+6\sqrt{2})^{2}/75+y_1^{2}/25=1\)．后者乘75：\((x_1+6\sqrt{2})^{2}+3y_1^{2}=75\)；由前者\(3y_1^{2}=3-x_1^{2}\)，代入：\(x_1^{2}+12\sqrt{2}x_1+72+3-x_1^{2}=75\)，即\(12\sqrt{2}x_1=0\)，\(x_1=0\)，\(y_1^{2}=1\)，\(y_1=\pm 1\)．\(A(0,\pm 1)\)．验算：此时\(B=(6\sqrt{2}/5,\pm 1/5)\)：\((72/25)/3+1/25=24/25+1/25=1\)成立，\(B\)确在椭圆上；\(|F_1A|=\sqrt{2+1}=\sqrt{3}\)，\(6\sqrt{2}/5-\sqrt{2}=\sqrt{2}/5\)，\(|F_2B|=\sqrt{2/25+1/25}=\sqrt{3}/5\)，比值\(\sqrt{3}/(\sqrt{3}/5)=5\)成立．}
\end{ansblock}}

\xiaojie{①识别：焦点三角形（周长、面积、张角存在性）与焦半径向量约束问题，判归本题型.}
\bindp {\kaishu ②操作：一切回到定义——周长拼\(4a\)、焦半径和\(2a\)；垂直即张角\(90^{\circ}\)，转勾股＋均值；向量等式设坐标代入方程组消元.}
\bindp {\kaishu ③收束：解出基本量后验存在性（判别式/取等点在不在椭圆上）；与\(B\)位置无关的定值题要说明理由而非举例.}

% ---------- 探究点四 综合构形与最值 ----------
\tjdnr{四}{综合构形与最值}{例\textbf{15}}{难(知识点三)}{已知椭圆\(C\)：\(x^{2}/4+y^{2}/3=1\)，\(F_1\)、\(F_2\)为\(C\)的左、右焦点，\(P\)为\(C\)上的一个动点（异于左右顶点），设\(\triangle F_1PF_2\)的外接圆面积为\(S_1\)，内切圆面积为\(S_2\)，则\(S_1+2S_2\)的最小值为\kongbai{}．}
\xiexwei{12mm}

{\ansblockgrayfalse
\begin{ansblock}[2章-练-课时11-难1]
% ans:2章-练-课时11-难1
\ansitem{15}{2π}
\ansnote{详解}{\(a=2\)、\(b=\sqrt{3}\)、\(c=1\)，\(|F_1F_2|=2\)．设\(\angle F_1PF_2=\theta\)．①\(\theta\)范围：\(P\)在短轴端点时\(\theta\)最大，此时\(|PF_1|=|PF_2|=2\)、\(|F_1F_2|=2\)为等边三角形，\(\theta_{\max}=\pi/3\)，故\(0<\theta\leqslant\pi/3\)．②外接圆：\(2R=|F_1F_2|/\sin\theta=2/\sin\theta\)，即\(R=1/\sin\theta\)．③内切圆：\(S_{\triangle F_1PF_2}=\frac{1}{2}|PF_1||PF_2|\sin\theta=b^{2}\tan(\theta/2)=3\tan(\theta/2)\)；又\(r=2S/\)周长\(=2\cdot 3\tan(\theta/2)/(2a+2c)=6\tan(\theta/2)/6=\tan(\theta/2)\)．④\(S_1+2S_2=\pi(R^{2}+2r^{2})=\pi[1/\sin^{2}\theta+2\tan^{2}(\theta/2)]\)．记\(t=\tan(\theta/2)\in(0,\sqrt{3}/3]\)（\(\theta\in(0,\pi/3]\)），\(\sin\theta=2t/(1+t^{2})\)：\(1/\sin^{2}\theta=(1+t^{2})^{2}/(4t^{2})\)，\(S_1+2S_2=\pi[(1+t^{2})^{2}/(4t^{2})+2t^{2}]=\pi[1/(4t^{2})+1/2+9t^{2}/4]\geqslant\pi\left[1/2+2\sqrt{\frac{1}{4t^{2}}\cdot\frac{9t^{2}}{4}}\right]=\pi(1/2+3/2)=2\pi\)，等号当\(1/(4t^{2})=9t^{2}/4\)，即\(t^{2}=1/3\)，\(t=\sqrt{3}/3\)为上限\(\theta=\pi/3\)（\(P\)为短轴端点）可取．最小值\(2\pi\)．}
\end{ansblock}}

\liB{变式\textbf{16}}{难(知识点三)}{}{已知椭圆\(C\)：\(x^{2}/a^{2}+y^{2}/b^{2}=1\)（\(a>b>0\)）的左、右焦点分别是\(F_1\)、\(F_2\)，\(P\)是椭圆上的动点，\(I\)和\(G\)分别是\(\triangle PF_1F_2\)的内心和重心，若\(IG\)与\(x\)轴平行，则椭圆的离心率为\nobreak(\kongwei)}
\bindopt {\fontsize{10.5pt}{\qpTJgLeadLiti}\selectfont \makebox[\dimexpr0.5\linewidth-1em\relax][l]{A．\(1/2\)}\makebox[\dimexpr0.5\linewidth-1em\relax][l]{B．\(\sqrt{3}/3\)}\par}

\bindopt {\fontsize{10.5pt}{\qpTJgLeadLiti}\selectfont \makebox[\dimexpr0.5\linewidth-1em\relax][l]{C．\(\sqrt{3}/2\)}\makebox[\dimexpr0.5\linewidth-1em\relax][l]{D．\(\sqrt{6}/3\)}\par}

{\ansblockgrayfalse
\begin{ansblock}[2章-练-课时11-难2]
% ans:2章-练-课时11-难2
\ansitem{16}{A}
\ansnote{详解}{\(O\)为\(F_1F_2\)中点，重心\(G\)在中线\(PO\)上且\(PG:GO=2:1\)．延长\(PI\)交\(x\)轴于\(Q\)（\(IQ\)在\(x\)轴内）．\(IG\parallel x\)轴，则在\(\triangle PQO\)中\(PI:IQ=PG:GO=2:1\)，即\(PQ:IQ=3\)．两三角形\(\triangle PF_1F_2\)与\(\triangle IF_1F_2\)同底\(F_1F_2\)，面积比\(=PQ:IQ=3\)．用内切圆半径\(r\)表示：\(S_{\triangle PF_1F_2}=\frac{1}{2}r(|PF_1|+|PF_2|+|F_1F_2|)\)（内心分三块）、\(S_{\triangle IF_1F_2}=\frac{1}{2}r|F_1F_2|\)，比\(=(|PF_1|+|PF_2|+|F_1F_2|)/|F_1F_2|=3\)，即\((|PF_1|+|PF_2|)/|F_1F_2|=2\)，即\(2a/2c=2\)，\(c/a=1/2\)．离心率为\(1/2\)，故选：A．}
\end{ansblock}}

\xiaojie{①识别：焦点三角形叠外接/内切圆、内心重心等特殊点构形的综合题，判归本题型.}
\bindp {\kaishu ②操作：设张角\(\theta\)统一表达各量；外接圆用正弦定理、内切圆用\(r=2S/\)周长；构形问题先用重心/内心比例把“平行”翻译成线段比，再转面积比.}
\bindp {\kaishu ③收束：换元后写明参数范围并核取等点在范围内（端点可取性）；求出的是比值还是离心率要回头看设问.}

% ============================================================
% 课堂评价 G5（恒5＝3单选+2填空＝导槽 G1~G5；ketang 新制顺排可跨栏，禁 \vbox 装栏）
% 印面号 17—21；多选门：本片正文多选 1（简6）≤4 合规
% ============================================================
\begin{ketang}{本组共5题（第17—21题），17—19为单选，20—21为填空．}

\jiancestem{{\fontsize{11.4pt}{14pt}\selectfont\heihao {\numboldjian 17}．}\hspace{4.1mm}\tieside{简单(知识点一)}\hspace{2.2mm}下列说法中正确的是\nobreak(\kongwei)}

\lxopt{A．已知\(F_1(-4,0)\)、\(F_2(4,0)\)，平面内到\(F_1\)、\(F_2\)两点的距离之和等于8的点的轨迹是椭圆}

\lxopt{B．已知\(F_1(-4,0)\)、\(F_2(4,0)\)，平面内到\(F_1\)、\(F_2\)两点的距离之和等于6的点的轨迹是椭圆}

\lxopt{C．平面内到两点\(F_1(-4,0)\)、\(F_2(4,0)\)的距离之和等于点\(M(5,3)\)到\(F_1\)、\(F_2\)的距离之和的点的轨迹是椭圆}

\lxopt{D．平面内到点\(F_1(-4,0)\)、\(F_2(4,0)\)距离相等的点的轨迹是椭圆}

\begin{ansblock}[2章-导-课时11-G1]
% ans:2章-导-课时11-G1
\ansitem{17}{C}
\ansnote{详解}{A：\(|F_1F_2|=8\)，距离和\(=8=|F_1F_2|\)，轨迹是线段\(F_1F_2\)，错；B：和\(6<8=|F_1F_2|\)，满足条件的点不存在，错；C：\(M(5,3)\)到两定点距离和\(=\sqrt{9^{2}+3^{2}}+\sqrt{1^{2}+3^{2}}=3\sqrt{10}+\sqrt{10}=4\sqrt{10}>8\)，和为常数且大于\(|F_1F_2|\)，轨迹是椭圆，对；D：到两点距离相等的点轨迹是线段\(F_1F_2\)的垂直平分线（\(y\)轴），错．故选：C．}
\end{ansblock}

\jiancestem{{\fontsize{11.4pt}{14pt}\selectfont\heihao {\numboldjian 18}．}\hspace{4.1mm}\tieside{简单(知识点一)}\hspace{2.2mm}已知\(F_1\)，\(F_2\)是两个定点，且\(|F_1F_2|=2a\)（\(a\)是正常数），动点\(P\)满足\(|PF_1|+|PF_2|=a^{2}+1\)，则动点\(P\)的轨迹是\nobreak(\kongwei)}

\bindopt {\fontsize{10.5pt}{\qpTJgLeadPingjia}\selectfont \makebox[\dimexpr0.5\linewidth-1em\relax][l]{A．椭圆}\makebox[\dimexpr0.5\linewidth-1em\relax][l]{B．线段}\par}

\bindopt {\fontsize{10.5pt}{\qpTJgLeadPingjia}\selectfont \makebox[\dimexpr0.5\linewidth-1em\relax][l]{C．椭圆或线段}\makebox[\dimexpr0.5\linewidth-1em\relax][l]{D．直线}\par}

\begin{ansblock}[2章-导-课时11-G2]
% ans:2章-导-课时11-G2
\ansitem{18}{C}
\ansnote{详解}{比较定值\(a^{2}+1\)与焦距\(2a\)：\(a^{2}+1-2a=(a-1)^{2}\geqslant 0\)，故\(a^{2}+1\geqslant 2a\)恒成立．当\(a>0\)且\(a\neq 1\)时\((a-1)^{2}>0\)，\(|PF_1|+|PF_2|>|F_1F_2|\)，轨迹为椭圆；当\(a=1\)时等号成立，\(|PF_1|+|PF_2|=|F_1F_2|\)，轨迹为线段\(F_1F_2\)．合起来：椭圆或线段，故选：C．}
\end{ansblock}

\jiancestem{{\fontsize{11.4pt}{14pt}\selectfont\heihao {\numboldjian 19}．}\hspace{4.1mm}\tieside{简单(知识点三)}\hspace{2.2mm}设\(F_1\)、\(F_2\)为椭圆\(x^{2}/25+y^{2}/16=1\)的两个焦点，直线过\(F_1\)交椭圆于\(A\)、\(B\)两点，则\(\triangle AF_2B\)的周长是\nobreak(\kongwei)}

\bindopt {\fontsize{10.5pt}{\qpTJgLeadPingjia}\selectfont \makebox[\dimexpr0.25\linewidth-0.5em\relax][l]{A．\(10\)}\makebox[\dimexpr0.25\linewidth-0.5em\relax][l]{B．\(15\)}\makebox[\dimexpr0.25\linewidth-0.5em\relax][l]{C．\(20\)}\makebox[\dimexpr0.25\linewidth-0.5em\relax][l]{D．\(25\)}\par}

\begin{ansblock}[2章-导-课时11-G3]
% ans:2章-导-课时11-G3
\ansitem{19}{C}
\ansnote{详解}{由定义\(|AF_1|+|AF_2|=2a=10\)，\(|BF_1|+|BF_2|=10\)；\(A\)、\(F_1\)、\(B\)共线故\(|AB|=|AF_1|+|BF_1|\)．周长\(=|AB|+|AF_2|+|BF_2|=|AF_1|+|BF_1|+|AF_2|+|BF_2|=20\)（\(=4a\)，与直线位置无关）．故选：C．}
\end{ansblock}

\jiancestem{{\fontsize{11.4pt}{14pt}\selectfont\heihao {\numboldjian 20}．}\hspace{4.1mm}\tieside{简单(知识点三)}\hspace{2.2mm}若椭圆\(x^{2}/25+y^{2}/9=1\)上一点\(P\)到两焦点的距离之积为\(m\)，则\(m\)的最大值为\kongbai{}．}
\xiexwei{10mm}

\begin{ansblock}[2章-导-课时11-G4]
% ans:2章-导-课时11-G4
\ansitem{20}{25}
\ansnote{详解}{\(a=5\)，由定义\(|PF_1|+|PF_2|=2a=10\)．均值不等式：\(m=|PF_1|\cdot|PF_2|\leqslant((|PF_1|+|PF_2|)/2)^{2}=25\)，等号当\(|PF_1|=|PF_2|=5\)．存在性验证：短轴端点\(P(0,\pm 3)\)到焦点\((\pm 4,0)\)的距离\(=\sqrt{16+9}=5\)，两边皆5，等号可取．\(m\)最大值为25．}
\end{ansblock}

\jiancestem{{\fontsize{11.4pt}{14pt}\selectfont\heihao {\numboldjian 21}．}\hspace{4.1mm}\tieside{中档(知识点三)}\hspace{2.2mm}椭圆\(x^{2}/12+y^{2}/3=1\)的一个焦点为\(F_1\)，点\(P\)在椭圆上，若线段\(PF_1\)的中点\(M\)在\(y\)轴上，则点\(M\)的纵坐标为\kongbai{}．}
\xiexwei{10mm}

\begin{ansblock}[2章-导-课时11-G5]
% ans:2章-导-课时11-G5
\ansitem{21}{±√3/4}
\ansnote{详解}{\(c^{2}=12-3=9\)，焦点\((\pm 3,0)\)．设另一定点为\(F_2\)．\(M\)为\(PF_1\)中点、\(O\)为\(F_1F_2\)中点，则\(OM\)为\(\triangle PF_1F_2\)的中位线，\(OM\parallel PF_2\)且\(OM\)在\(y\)轴上，故\(PF_2\perp x\)轴，\(P\)横坐标为\(\pm 3\)．代入：\(9/12+y^{2}/3=1\)，即\(y^{2}=3/4\)，\(y_P=\pm\sqrt{3}/2\)．\(M\)纵坐标\(=y_P/2=\pm\sqrt{3}/4\)．}
\end{ansblock}

\end{ketang}

% ---- 尾块（\tailfill 置 \end{multicols} 前；无参＝内容尾随框，件侧不擅自定高） ----
\tailfill

\end{multicols}
\end{document}
